% ============================================================================
% IEEE Transactions on Transportation Electrification Paper Structure
% LLAMBO-MO: LLM-Augmented Multi-Objective Bayesian Optimization
% for Battery Fast-Charging Protocol Design
% ============================================================================

\documentclass[journal]{IEEEtran}
\usepackage{cite}
\usepackage{amsmath,amssymb,amsfonts}
\usepackage{algorithm}
\usepackage{algpseudocode}
\usepackage{graphicx}
\usepackage{textcomp}
\usepackage{xcolor}
\usepackage{booktabs}
\usepackage{multirow}
\usepackage{subfigure}

\hyphenation{op-tical net-works semi-conduc-tor}

\begin{document}

% ============================================================================
% Title
% ============================================================================
\title{
LLAMBO-MO: Large Language Model-Augmented Multi-Objective Bayesian Optimization \\
for Battery Fast-Charging Protocol Design
}

\author{
First~A.~Author,~\IEEEmembership{Member,~IEEE,}
Second~B.~Author,~\IEEEmembership{Fellow,~OSA,}
and~Third~C.~Author,~\IEEEmembership{Life~Fellow,~IEEE}%
\thanks{Manuscript received Month Day, Year; revised Month Day, Year.
This work was supported in part by ...}
\thanks{First A. Author is with the ...}
\thanks{Second B. Author is with the ...}
\thanks{Third C. Author is with the ...}
}

\markboth{IEEE TRANSACTIONS ON TRANSPORTATION ELECTRIFICATION,~Vol.~XX, No.~XX, Month~2026}%
{Author \MakeLowercase{\textit{et al.}}: LLAMBO-MO: LLM-Augmented Multi-Objective BO for Battery Fast-Charging}

\maketitle

% ============================================================================
% Abstract (150-250 words)
% ============================================================================
\begin{abstract}
The design of fast-charging protocols for lithium-ion batteries is inherently
a constrained multi-objective optimization problem (CMOP) that must simultaneously
minimize charging time, peak temperature rise, and capacity fade.
While Bayesian optimization (BO) has emerged as a sample-efficient paradigm
for such expensive black-box problems, existing methods rely solely on
mathematical surrogates and fail to leverage the rich domain knowledge
embedded in battery literature.
This paper proposes LLAMBO-MO, a large language model (LLM)-augmented
multi-objective BO framework that strategically integrates LLM-generated
domain knowledge at two touchpoints within the BO loop.
First, the LLM generates physics-informed warmstart candidates to bootstrap
the surrogate with high-quality initial data.
Second, the LLM provides iteration-level guidance that is coupled with
the Gaussian process posterior through a bounded acquisition-time mean shift.
The multi-objective problem is decomposed via augmented Tchebycheff
scalarization with Riesz s-energy weight vectors.
Experiments on an SPMe-based electrochemical-thermal-aging simulator
demonstrate that LLAMBO-MO outperforms ParEGO, NSGA-II, DISK, and PIMD
in both hypervolume convergence and Pareto front quality within a fixed
evaluation budget.
\end{abstract}

\begin{IEEEkeywords}
Lithium-ion battery, fast charging, multi-objective optimization,
Bayesian optimization, large language models, Gaussian process,
Tchebycheff scalarization.
\end{IEEEkeywords}

\IEEEpeerreviewmaketitle

% ============================================================================
% Section I: Introduction
% ============================================================================
\section{Introduction}
\label{sec:introduction}

% Recommended structure (3-4 pages):
% 1. Background: Battery fast charging importance and challenges
% 2. Problem formulation: CMOP with three competing objectives
% 3. Existing approaches: ECM, ROM, BO-based methods
% 4. Gap: BO methods ignore domain knowledge; LLM potential untapped
% 5. Contributions: Three main contributions
% 6. Organization: Paper structure overview

% Key points to cover:
% - Battery electrification context [1]-[3]
% - Fast charging trade-off: speed vs. degradation [4], [5]
% - Model-based optimization advantages [6]-[8]
% - BO for charging protocol design [4], [17], [18]
% - LLM capabilities for scientific knowledge [Ref-LLM]
% - LLAMBO-MO motivation and novelty

% ============================================================================
% Section II: Problem Formulation
% ============================================================================
\section{Problem Formulation}
\label{sec:problem}

% 2.1 Multi-Objective Optimization Problem
\subsection{Charging Protocol Design as CMOP}

% Objectives:
% min f(theta) = {t_c, Delta T_p, Q_s}
% where t_c: charging time
%       Delta T_p: peak temperature rise
%       Q_s: capacity fade

% 2.2 Decision Variables
\subsection{Decision Variables and Constraints}

% Variables: theta = [I_1, I_2, I_3, Delta s_1, Delta s_2]
% Bounds and constraint table

% 2.3 Operational Constraints
\subsection{Operational Constraints}

% Voltage, temperature, SOC limits during simulation

% ============================================================================
% Section III: Electrochemical-Thermal-Aging Model
% ============================================================================
\section{Electrochemical-Thermal-Aging Model}
\label{sec:model}

% 3.1 Electrochemical Model
\subsection{Electrochemical Model}

% SPMe model description [23], [24]
% PyBaMM implementation [PyBaMM-ref]

% 3.2 Thermal Model
\subsection{Thermal Model}

% Lumped-parameter thermal model
% Heat generation equation

% 3.3 Aging Model
\subsection{Aging Model}

% Empirical model (Eq. 6-7)
% Physics-based SEI growth model (Eq. 8)

% 3.4 Simulation Implementation
\subsection{Simulation Implementation}

% Stage decomposition
% Stage-wise solving
% Objective extraction

% ============================================================================
% Section IV: Proposed LLAMBO-MO Framework
% ============================================================================
\section{Proposed LLAMBO-MO Framework}
\label{sec:method}

% 4.1 Framework Overview
\subsection{Motivation and Framework Overview}

% Two-touchpoint architecture diagram
% Key insight: LLM complements GP at strategic points

% 4.2 Decomposition Strategy
\subsection{Riesz s-Energy Weight Generation and Tchebycheff Scalarization}

% Log-transform and normalization (Eq. 9-10)
% Augmented Tchebycheff scalarization (Eq. 11)
% Riesz s-energy relaxation (Eq. 12-13)

% 4.3 GP Surrogate Model
\subsection{Gaussian Process Surrogate Model}

% GP posterior (Eq. 14)
% Matern 5/2 kernel with ARD (Eq. 15)
% Hyperparameter optimization

% 4.4 LLM Integration (Two Touchpoints)
\subsection{LLM Integration: Two-Touchpoint Architecture}

\subsubsection{Touchpoint 1: Warmstart Candidate Generation}
% Prompt design
% Portfolio selection
% Fallback mechanisms

\subsubsection{Touchpoint 2: Iteration-Level Guidance}
% Context provided to LLM
% Point mode vs. Region mode

% 4.5 GP-LLM Coupling
\subsection{GP-LLM Coupling via Acquisition-Time Mean Shift}

% Grid construction and weighting
% Coupling strength lambda (Eq. 16)
% Acquisition-time mean shift (Eq. 17-18)

% 4.6 Acquisition Function
\subsection{Acquisition Function with Stagnation-Aware Exploration}

% Expected Improvement (Eq. 19)
% Acquisition prior (Eq. 20)
% Stagnation detection (Eq. 21)

% 4.7 Overall Algorithm
\subsection{Overall Algorithm}

% Algorithm 1: LLAMBO-MO pseudocode

% ============================================================================
% Section V: Experiments
% ============================================================================
\section{Experiments}
\label{sec:experiments}

% 5.1 Experimental Setup
\subsection{Experimental Setup}

% Battery models: Chen2020, Ecker2015
% Decision variables and objectives
% Evaluation budget: 56 evaluations
% Performance metric: normalized HV

% 5.2 Baseline Methods
\subsection{Baseline Methods}

% ParEGO [Ref-ParEGO]
% NSGA-II [Ref-NSGA2]
% DISK [Ref-DISK]
% PIMD [Ref-PIMD]

% 5.3 HV Convergence Comparison
\subsection{HV Convergence Comparison}

% Table III: Final HV on Chen2020
% Fig. X: HV convergence curves
% Table IV: Results on Ecker2015

% 5.4 Pareto Front Quality
\subsection{Pareto Front Quality}

% Fig. X: Pareto front visualization
% Representative protocols at extremes and middle

% 5.5 Ablation Study
\subsection{Ablation Study (V0--V6)}

% Table V: Ablation results
% Analysis of Warmstart vs. Iteration guidance contribution

% 5.6 Computational Efficiency
\subsection{Computational Efficiency}

% Runtime comparison
% Overhead breakdown

% ============================================================================
% Section VI: Conclusion
% ============================================================================
\section{Conclusion}
\label{sec:conclusion}

% Summary of contributions
% Experimental validation
% Future work directions

% ============================================================================
% References
% ============================================================================
\begin{thebibliography}{1}

\bibitem{ref1}
N.~Nitta \textit{et al.}, ``Li-ion battery materials: present and future,''
\textit{Materials Today}, vol.~18, no.~5, pp.~252--264, 2015.

\bibitem{ref2}
A.~Tomaszewska \textit{et al.}, ``Lithium-ion battery fast charging: A review,''
\textit{eTransportation}, vol.~1, p.~100011, 2019.

% ... (remaining references)

\end{thebibliography}

\end{document}
